% ----------------------------------------------------------------------
%  Введение к главе 3
% ----------------------------------------------------------------------
В главах~1 и~2 нас интересовала геометрия: как расположены атомы и как
нарушается их порядок. Эта глава отвечает на вопрос \emph{почему}:
почему атомы вообще собираются в кристалл, почему одни вещества~—
металлы, а другие~— изоляторы, почему металл проводит ток и тепло,
откуда берутся ферромагнетизм и тепловое расширение. Ответы почти
всегда сводятся к поведению электронов и к колебаниям атомов решётки.

\begin{figure}[ht]
\centering
\begin{tikzpicture}[
  box/.style={draw=#1, fill=#1!8, rounded corners=2pt, align=center,
    font=\small, inner sep=4pt, text width=40mm, minimum height=11mm, line width=0.6pt},
  flow/.style={-{Stealth[length=2mm]}, line width=0.6pt, fmGray}]
  \node[box=fmBlue] (a) at (0,0) {\textbf{3.1} Атомы и типы\\межатомной связи};
  \node[box=fmBlue] (b) at (0,-1.75) {\textbf{3.2} Электроны в кристалле:\\свободные электроны, зоны};
  \node[box=fmGreen] (c) at (-4.6,-3.6) {\textbf{3.3} Фазы в сплавах};
  \node[box=fmGreen] (d) at (-4.6,-5.3) {\textbf{3.4} Диффузия и\\фазовые превращения};
  \node[box=fmAmber] (e) at (0,-3.6) {\textbf{3.5} Электрические\\свойства};
  \node[box=fmAmber] (f) at (0,-5.3) {\textbf{3.6} Магнитные\\свойства};
  \node[box=fmRed] (g) at (4.6,-3.6) {\textbf{3.7} Тепловые свойства\\(фононы)};
  \node[box=fmRed] (h) at (4.6,-5.3) {\textbf{3.8} Упругие свойства,\\внутреннее трение};
  \draw[flow] (a) -- (b);
  \draw[flow] (b) -- (e); \draw[flow] (b.west) -| (c.north);
  \draw[flow] (b.east) -| (g.north);
  \draw[flow] (e) -- (f); \draw[flow] (c) -- (d); \draw[flow] (g) -- (h);
\end{tikzpicture}
\caption{Карта главы. Из природы межатомной связи и поведения электронов
вытекают строение сплавов (зелёный блок), электрические и магнитные
(жёлтый) и тепловые и упругие (красный) свойства}
\label{fig:ch3-map}
\end{figure}
